\documentclass[a4paper,twoside]{article}


% --- RUIMTE EN LAYOUT ---
\usepackage{setspace}

\usepackage{geometry} %size regulating
\geometry{
 a4paper,
 left=25mm,
 right=25mm,
 top=20mm,
 bottom=20mm,
 }

\usepackage{parskip} % for better spacing between paragraphs
%\usepackage{flushend} % Alleen nodig als je laatste pagina van artikel wilt vullen met tekst, zodat de kolommen even lang zijn. Anders kan je dit beter weglaten.

% --- WISKUNDE ENSYMBOLEN ---
\usepackage{amsmath, amssymb, amsfonts, amsthm} % For nicer math format
\usepackage{mathrsfs} % For curly arrr
\usepackage{bm} % For bold symbols eg \bm{\nabla}
\usepackage{esint} % For bolletje in integrals
\usepackage{dsfont}
\usepackage{siunitx} % Voeg deze expliciet toe voor die graden en eenheden!
\usepackage{braket} % For bra-ket notation
\usepackage{mathtools}
\usepackage{url}
\allowdisplaybreaks

% --- OVERIGE ---

\usepackage{graphicx} % Required for inserting images
\usepackage[hidelinks]{hyperref} % geen kaders om linken
\usepackage{subcaption}
%\usepackage[dutch]{babel}
\usepackage[T1]{fontenc}
\usepackage[utf8]{inputenc}
\usepackage{lmodern}
\usepackage{multicol}


% --- Tikz packages ---
\usepackage{tikz}
\usetikzlibrary{angles,quotes, babel} % for pic (angle labels)
\usetikzlibrary{arrows.meta,decorations.markings,calc, decorations.pathmorphing, decorations.pathreplacing, positioning, shapes, shadings, patterns}
\usepackage{xcolor}
\usepackage{tikz-3dplot}
\usepackage[siunitx]{circuitikz}

\usepackage{etoolbox} % ifthen
\usepackage[outline]{contour} % glow around text
\usepackage{xfp} % higher precision (16 digits?)
\contourlength{1.1pt}



% --- THEOREM STIJLEN MET EXTRA RUIMTE ---

\theoremstyle{definition}
\newtheorem{postulate}{Postulaat}
\newtheorem{definition}{Definitie}
\newtheorem{derivation}{Afleiding}
\newtheorem{proposition}{Propositie}
\newtheorem{example}{Voorbeeld}

\theoremstyle{plain}
\newtheorem{theorem}{Theorem}


% --- EXTRA COMMANDO'S ---

\newcommand{\dbar}{{\mkern+3mu\mathchar'26\mkern-12mu d}}

% --- Titel en Introductie --- 


\title{Statistische Fysica}
\author{T. Cools}
\date{2026-2027}

\begin{document}

\maketitle

\begin{abstract}
    Dit is mijn versie van de cursus Statistische Fysica, Waar ik mezelf test op wat ik weet en begrijp van de theorie.
     De inhoud is gebaseerd op de cursus van Prof. Dr. J. Ryckebusch.
\end{abstract}

\tableofcontents


\section{Eerste Wet van de Thermodynamica}
\subsection{Enkele thermische begrippen}
\subsection{De eerste wet van de thermodynamica}
\subsection{Reversibele en irreversibele processen}

\section{Tweede Wet van de Thermodynamica}
\subsection{De tijdsevolutie van natuurlijke processen}
\subsection{Het statistisch gewicht van een macrotoestand}
\subsection{Evenwicht van een ge\"isoleerd systeem}
\subsection{Schottky defecten}
\subsection{Equipartitietheorema en het viriaal theorema}
\subsection{Theorema van Liouville}
\subsection{Evenwicht van een systeem in een warmtebad}
\subsection{Onzekerheid, informatie, entropie en de Boltzmann distributie}

\section{Paramagnetisme}
\subsection{Paramagnetische vaste stof in een warmtebad}
\(\bm{B}\) is het extern magnetisch veld \(\bm{H}\) is ook een extern veld met de relatie \(\bm{B} \equiv \mu_0 \bm{H}\) waarbij \(\mu_0\) de magnetische permitiviteit van het vacuum is.

Energie:
\begin{itemize}
    \item parallel aligned: \(E_+ = - \mu \bm B\)
    \item anti-parallel aligned: \(E_- = \mu \bm B\)
\end{itemize}


De partitiefunctie: 
\begin{align*}
    Z(T,\bm B, N) 
    &= Z(T,\bm B, N=1)\\ 
    &= e^{\frac{E_+}{kT}}+e^{\frac{E_-}{kT}}
    = e^{\frac{- \mu \bm B}{kT}}+e^{\frac{\mu \bm B}{kT}}\\
    &= e^{\frac{\mu \bm B}{kT}}+e^{\frac{-\mu \bm B}{kT}}\\
    &= 2 \cosh(\frac{\mu \bm B}{kT}) =  2 \cosh x
\end{align*}

Merk op \(x \equiv \frac{\mu \bm B}{kT}\) is dimensieloos!

Waarschijnlijkheid \(p_\pm\):
\begin{align*}
        p_\pm 
        &= \frac{e^{\pm \frac{\mu \bm B}{kT}}}{Z(T,\bm B, N)}\\
        &= \frac{e^{\pm x}}{2 \cosh x}
\end{align*}


Dan hebben we ook nog de average values ingevoerd:

\(\overline{\mu} = \mu (p_+ - p_-)= \mu \tanh x \equiv  kT \left(\frac{\partial \ln Z_1}{\partial \bm B}\right)_T = \frac{1}{\beta} \left(\frac{\partial \ln Z_1}{\partial \bm B}\right)_T\)

en
\(\overline{E_1} = -\mu \bm B (p_+ - p_-) = -\mu \bm B \tanh x \equiv -\left(\frac{\partial \ln Z_1}{\partial \beta}\right)_B\)

\(\overline{E} = N \overline{E_1} = -\mu \bm B N \tanh x \)

\(\bm M = N \overline{\mu} = \mu N \tanh x \)

\subsection{Warmtecapaciteit en entropie}

\subsection{Absolute negatieve temperatuur}

\section{Evenwichtscondities van Systemen}
\subsection{De eerste wet voor infinitesimale veranderingen}
\subsection{De ongelijkheid van Clausius}
\subsection{Eigenschappen van de Helmholtz vrije energie}
\subsection{Andere thermodynamische potentialen}
\subsection{De derde wet van de thermodynamica}

\section{Warmtecapaciteit van Vaste Stoffen}
\subsection{Inleiding}
\subsection{Het model van Einstein}
\subsection{Het model van Debye}

\section{Klassiek Gas}
\subsection{De definitie van een klassiek en ideaal gas}
\subsection{De partitiefunctie voor een klassiek en ideaal gas}
\subsection{Geldigheidscriterium voor het klassiek regime}
\subsection{De toestandsvergelijking}
\subsection{De entropie van een klassiek en ideaal gas}
\subsection{De Maxwell snelheidsdistributiefunctie}
\subsection{Re\"ele gassen}

\section{Vloeistoffen en Moleculaire Dynamica}
\subsection{Inleidende opmerkingen}
\subsection{Paar correlatiefunctie}
\subsubsection{Theorie}
\subsubsection{Experimentele bepaling}
\subsubsection{Verband met observabelen}
\subsection{Moleculaire Dynamica}
\subsubsection{Beschrijving van de techniek}
\subsubsection{Een concreet voorbeeld}
\subsection{Toepassingen van Moleculaire Dynamica}
\subsubsection{Random walk en diffusie}
\subsubsection{Berekening van radiale distributiefuncties}
\subsubsection{Verband tussen diffusie en snelheidsfluctuaties}

\section{Kwantumgassen}
\subsection{Inleidende opmerkingen}
\subsection{Statistiek voor ideaal kwantumgas in het $(T,V,N)$ systeem}
\subsection{De partitiefunctie van een ideaal kwantumgas}
\subsection{Kwantum microtoestanden en kwantum macrotoestanden}
\subsubsection{Dichtheidsmatrixformalisme}
\subsubsection{Tijdsafhankelijkheid van de kwantum dichtheidsmatrix}

\section{Ideaal Fotonengas}
\subsection{Inleidende opmerkingen}
\subsection{De partitiefunctie voor fotonen}
\subsection{De Wet van Planck}
\subsection{Eigenschappen van een Zwarte Straler}

\section{Groot-Canonisch Systeem}
\subsection{De Gibbs distributie}
\subsubsection{Inleiding}
\subsubsection{Afleiding van de Gibbs of Groot-Canonische distributiefunctie}
\subsection{De FD en BE distributie voor een ideaal kwantumgas}
\subsubsection{Gemiddeld bezettingsgetal in een ideaal kwantumgas}
\subsubsection{Fluctuaties in de bezettingsgetallen in een ideaal kwantumgas}
\subsection{De thermodynamica van een Gibbs systeem}
\subsection{Toestandsvergelijking voor ideale kwantumgassen}
\subsection{Het Groot-Canonisch systeem en re\"ele gassen}

\section{Ideaal Fermi Gas}
\subsection{De FD energiedistributie}
\subsection{De toestandsvergelijking van een ideaal Fermi gas}
\subsubsection{Algemene uitdrukking}
\subsubsection{De functies $f_{3/2}(z)$ en $f_{5/2}(z)$ bij kleine en grote $z$}
\subsubsection{Kleine $z$: toestandsvergelijking in de klassieke limiet}
\subsubsection{Grote $z$: toestandsvergelijking in het kwantummechanisch regime}
\subsection{Het aandeel van elektronen in de soortelijke warmte van metalen}
\subsection{Pauli paramagnetisme}
\subsection{Het Relativistisch Fermi Gas: Witte Dwergen}

\section{Ideaal Bose Gas}
\subsection{De BE energiedistributie}
\subsection{Toestandsvergelijking van een ideaal Bose gas}
\subsection{Soortelijke warmte en entropie van een ideaal Bose gas}
\subsection{Bose-Einstein condensatie in dunne atomaire gassen}
\subsection{Supervloeibaarheid}

\section{Theorie van Tweede-orde Faseovergangen}
\subsection{Inleiding}
\subsection{Het Ising systeem}
\subsection{Het eendimensionaal Ising systeem}
\subsubsection{Exacte oplossing}
\subsubsection{Spin-spin correlatiefunctie}
\subsection{Het tweedimensionaal Ising Systeem}
\subsubsection{Computersimulaties}
\subsubsection{Exacte oplossing}
\subsection{Gemiddeld-veld theorie van het Ising systeem}
\subsection{Universaliteit}
\subsection{Het rooster gas}
\subsection{Orde/wanorde transities in legeringen en universaliteit}

\appendix

\section{Toestandsdichtheden}
\subsection{Aantal staande-golf oplossingen in een volume $V$}
\subsection{Elastische golven in een continu medium}

\section{Gaussische Integralen, Volumes van Hypersferen}
\subsection{Gaussische Integralen}
\subsection{Berekening van het volume van een $n$-dimensionale bol met straal $R$}
\subsection{Dichtheid der toestanden voor $3N$ translationele vrijheidsgraden}

\end{document}